\subsection{Проверьте себя}
\begin{enumerate}
  \item Как начать миссию? \verb|ap.push(___)|
    \begin{enumerate}[label=(\alph*)] \item \verb|Ev.MCE_PREFLIGHT| \item \verb|Ev.START| \item \verb|Ev.MCE_TAKEOFF| \end{enumerate}
    Ответ: (a) --- предстарт всегда первый.
  \item Какое событие сообщает, что взлёт завершён?
    \begin{enumerate}[label=(\alph*)] \item \verb|Ev.ENGINES_STARTED| \item \verb|Ev.TAKEOFF_COMPLETE| \item \verb|Ev.POINT_REACHED| \end{enumerate}
    Ответ: (b).
  \item Как задать курс $90^\circ$?
    \begin{enumerate}[label=(\alph*)] \item \verb|ap.updateYaw(90)| \item \verb|ap.updateYaw(math.rad(90))| \item \verb|ap.updateYaw(90 / 180)| \end{enumerate}
    Ответ: (b) --- радианы.
  \item Когда отправлять вторую точку маршрута?
    \begin{enumerate}[label=(\alph*)] \item Через 4 с после первой \item По \verb|POINT_REACHED| первой \item Сразу вслед за первой \end{enumerate}
    Ответ: (b).
  \item Как безопасно выполнить действие через 2 секунды?
    \begin{enumerate}[label=(\alph*)] \item \verb|sleep(2); action()| \item \verb|Timer.callLater(2, action)| \item \verb|Timer.new(2, action)| \end{enumerate}
    Ответ: (b).
  \item Как организовать непрерывное обновление точки окружности?
    \begin{enumerate}[label=(\alph*)] \item \verb|while true do ap.goToLocalPoint(...) end| \item \verb|Timer.new(0.1, function() ap.goToLocalPoint(...) end)| \item \verb|repeat ... until false| \end{enumerate}
    Ответ: (b).
  \item Что делать по \verb|Ev.SHOCK|?
    \begin{enumerate}[label=(\alph*)] \item Продолжить маршрут \item \verb|ap.push(Ev.MCE_LANDING)| \item Остановить таймеры и \verb|ap.push(Ev.ENGINES_DISARM)| \end{enumerate}
    Ответ: (c).
  \item Дрон прилетел в точку $(2, 0, 1)$, \verb|Sensors.lpsPosition()| вернул \verb|x = 1.93|. Как правильно проверить прибытие?
    \begin{enumerate}[label=(\alph*)] \item \verb|x == 2| \item \verb|math.abs(2 - x) <= 0.15| \item \verb|x > 2| \end{enumerate}
    Ответ: (b) --- измерение неточно, сравнивают с допуском.
  \item Где в \verb|rc_chans = table.pack(Sensors.rc())| тумблер SwA?
    \begin{enumerate}[label=(\alph*)] \item \verb|rc_chans[1]| \item \verb|rc_chans[8]| \item \verb|rc_chans["SwA"]| \end{enumerate}
    Ответ: (b).
\end{enumerate}

\subsection{Практика: проекты миссий}

Решения и дополнительные идеи (спираль, лоитер, лепесток, радиус с пульта, ограничение скорости поворота, повтор после посадки) --- в \texttt{examples/lua-tasks/07\_advanced\_missions/}. Во всех проектах: предстарт и взлёт как в примерах, светодиодная индикация этапов, реакции на \verb|SHOCK| и \verb|LOW_VOLTAGE2|, а в журнале --- строки с подписанными этапами и причиной остановки. Проект принят, если всё перечисленное в «Что должно получиться» видно в сцене или в журнале.

\begin{taskbox}[Многоугольник из $N$ точек]
Сгенерируйте циклом $N$ вершин окружности радиуса 0.5~м под углами $360^\circ \cdot i / N$, пройдите их счётчиком шагов, вернитесь в центр и сядьте. Проверьте для $N = 3, 6, 12$, меняя только одно число.
\expected{след --- правильный треугольник, шестиугольник, двенадцатиугольник с центром в начале сцены; в журнале по строке на каждую достигнутую вершину с её номером; посадка в центре; при вызове \texttt{callback(Ev.LOW\_VOLTAGE2)} посреди маршрута дрон садится на месте, линейка красная.}
\end{taskbox}

\begin{taskbox}[Два квадрата]
Пролетите квадрат со стороной 1~м дважды подряд и только потом садитесь. Подсказка: счётчик кругов рядом со счётчиком шагов.
\expected{на сцене --- один квадрат, пройденный дважды (след ложится сам на себя); в журнале номера кругов 1 и 2 и восемь достигнутых вершин; посадка только после восьмой.}
\end{taskbox}

\begin{taskbox}[Восьмёрка с курсом]
Две окружности с противоположным направлением обхода и общей точкой в центре; нос всё время по касательной; завершение --- в центре и посадка. Окружности ведёт таймер, переходы --- конечный автомат с подсветкой состояний.
\expected{след --- восьмёрка; камера \ui{Дрон} показывает, что нос смотрит по направлению движения на обеих петлях; цвета линейки меняются при каждом переходе; посадка в центре, подтверждённая \texttt{near}.}
\end{taskbox}

\begin{taskbox}[Старт и останов тумблером]
Тумблер SwA (\verb|rc_chans[8]|) вверх --- предстарт, взлёт, окружность; тумблер вниз --- прекратить, вернуться в точку над стартом, сесть. Проверьте в панели \ui{Пульт}.
\expected{пока тумблер внизу, дрон стоит; после переключения вверх взлетает и кружит; после переключения вниз в любой момент окружности возвращается и садится; повторное переключение во время посадки ничего не ломает; в журнале --- причина остановки «тумблер».}
\end{taskbox}

\begin{taskbox}[Отчёт о точности]
Пройдите 8 точек окружности. По каждому \verb|POINT_REACHED| печатайте номер точки, время с предыдущего события (\verb|time()| возвращает секунды) и ошибку положения --- расстояние от измеренного положения до цели. В конце напечатайте среднюю и наибольшую ошибку.
\expected{восемь строк отчёта и итоговая строка; все ошибки меньше 0.15~м --- допуска, с которым автопилот засчитывает прибытие; интервалы между точками близки друг к другу. По отчёту видно, на какой точке ошибка была наибольшей.}
\mission{Lua}{72}
\end{taskbox}

\subsection*{Итоги главы}
\begin{itemize}
    \item Миссия --- цепочка «команда $\to$ событие $\to$ команда»; цепочку ведёт \verb|callback|.
    \item Предстарт, взлёт, точки, посадка; счётчик шагов для маршрута; таймер для плавных траекторий.
    \item Координаты в метрах от начала сцены, Z вверх; на старте нос смотрит вдоль +Y; курс в радианах.
    \item Сложная миссия --- конечный автомат: состояние проверяется вместе с событием.
    \item Failsafe обязателен: \verb|SHOCK| --- моторы выключить, \verb|LOW_VOLTAGE2| --- сесть.
    \item Результат проверяют измерением: \verb|Sensors.lpsPosition()|, сравнение с допуском, тайм-аут и причина остановки в журнале.
\end{itemize}
